% 第3章3.5.2交替路径翻转：允许边 a-x, a-y, b-y, b-z, c-y；当前匹配 {a-x, b-y}（已是极大匹配）。
% 从自由左顶点 c 出发：c-y（未匹配）、y-b（匹配）、b-z（未匹配），z 自由，得到增广路径；
% 翻转后匹配 {a-x, c-y, b-z}，大小由2增为3。
\begin{tikzpicture}[line width=.8pt,font=\fontsize{8}{10}\selectfont,
  v/.style={circle,draw=FigureStroke,fill=FigureBlueFill,line width=.8pt,minimum size=6.5mm,inner sep=0pt,text=FigureInk},
  fv/.style={v,fill=FigureYellowFill},
  e/.style={draw=FigureMuted,line width=.8pt},
  m/.style={draw=FigureBlue,line width=1.8pt},
  pu/.style={draw=FigureGreen,line width=1.5pt,dashed},
  pm/.style={draw=FigureRed,line width=1.8pt},
  note/.style={text=FigureMuted,font=\FigureNoteFont\fontsize{7.5}{9}\selectfont}]
  \foreach \dx/\k in {0/A,4.2/B}{
    \node[v] (a\k) at (\dx,1.1) {$a$};
    \node[v] (b\k) at (\dx,0) {$b$};
    \node[v] (c\k) at (\dx,-1.1) {$c$};
    \node[v] (x\k) at (\dx+2,1.1) {$x$};
    \node[v] (y\k) at (\dx+2,0) {$y$};
    \node[v] (z\k) at (\dx+2,-1.1) {$z$};
  }
  \node[fv] at (cA) {$c$}; \node[fv] at (zA) {$z$};
  % (a)
  \draw[m] (aA)--(xA);
  \draw[e] (aA)--(yA);
  \draw[pm] (bA)--(yA);
  \draw[pu] (bA)--(zA);
  \draw[pu] (cA)--(yA);
  \node at (1,-1.95) {(a) 匹配$\{a-x,b-y\}$与增广路径};
  % (b)
  \draw[m] (aB)--(xB);
  \draw[e] (aB)--(yB);
  \draw[e] (bB)--(yB);
  \draw[m] (bB)--(zB);
  \draw[m] (cB)--(yB);
  \node at (5.2,-1.95) {(b) 翻转后$\{a-x,b-z,c-y\}$};
  \node[note] at (3.1,-2.55) {粗蓝线：匹配边；红线：路径上将撤去的匹配边；绿色虚线：路径上将加入的边；黄色：自由端点};
\end{tikzpicture}
